\subsection*{接班者怎样知道旧决定}

前面用 Raft 说明怎样延续已提交前缀。现在换成 Paxos，保留同一个问题：旧处理者已经建立决定，响应却没有到达，新处理者从哪里得到继续依据？

先只决定日志中的一个位置，暂定为 41。参与者包括提出候选值的提议者、保存承诺与接受事实的接受者，以及获知决定的学习者。这些是协议角色，同一个服务进程可以兼任。候选值可以是完整 S1 购买命令。Paxos 用唯一且可比较的提议编号区分不同轮次，一份提议包含编号与值。下面把编号写为 7、8，假定分配机制保证全局唯一。编号标记这次争取决定的轮次，位置 41 标记最终要决定哪个命令槽~\cite{paxos-made-simple}。

设接受者为 A、B、C，任意两个形成多数。提议者 L1 已经按协议完成编号 7 的准备阶段，接着发出接受请求 \code{(7,S1)}。A、B 持久接受，C 尚未收到，L1 失效。位置 41 的值此时已经选定为 S1，即使没有学习者及时收到足够回信。客户端能否获知，取决于后来怎样收集并返回这些接受事实。\footnote{《Paxos Made Simple》指出：“a value could be chosen with no learner ever finding out”。值被选定和学习者获知是两个事实。工程实现可以经准备阶段及多数接受者保存的状态恢复~\cite{paxos-made-simple}。}

L2 自己的本地位置 41 为空时，先向接受者发送更高编号 8 的准备请求。接受者答复时，承诺今后不接受小于 8 的提议，并报告自己已经接受的最高编号提议。L2 从 B、C 收到答复：B 报告 \code{(7,S1)}，C 报告尚未接受。L2 取这些答复里编号最高的已接受值，把 S1 作为编号 8 的值，再进入接受阶段。获得多数接受后，它可以把 S1 的决定传给学习者~\cite{paxos-made-simple}。

\CElevenFigure{ch11-original16}{一个位置的 Paxos 恢复：准备阶段取回旧值并阻止更小编号继续形成新决定}{fig:ch11-paxos-one-slot}

这个准备阶段同时处理过去和未来。B 的答复把已经存在的 S1 带给新提议者，B/C 的承诺又挡住迟到的编号 7 接受请求。假如编号 7 在准备以前已经选定值，它的接受多数与准备多数必有交集，至少一份相关接受事实会进入答复。更高轮次按答复中最高已接受编号的值继续，归纳地把已选定值保留下来。

若准备多数全都报告没有接受过值，新提议者可以自由选择。因为任何较小编号的接受多数都会与这份准备多数相交，且交集成员既没有接受过旧值，又承诺阻止后来的较小编号请求，旧轮次便无法绕开它形成另一个决定。实际情况中，准备者也可能看到尚未达到多数的候选。按最高已接受值继续，是在当前证据下保持安全的规则。

接受者把最高承诺编号与最高已接受提议保存到稳定存储，并在答复以前完成所需持久化。它崩溃后若忘记承诺，就可能重新接受已被更高轮次排除的消息。若忘记接受值，新提议者又可能从空答复中得出错误选择。唯一提议编号、持久承诺、值的继承和多数交集共同支撑这里的论证~\cite{paxos-made-simple,part-time-parliament}。

\subsection*{从一个位置到连续状态}

购买服务要处理许多命令。多实例 Paxos 为位置 1、2、3 等分别决定值。通常由稳定领导者为客户端命令分配位置，并把多个位置的准备工作合并完成，后续正常命令主要进行接受阶段。这种工程组织常称为 Multi-Paxos~\cite{paxos-made-simple}。

不同位置的决定可能出现空洞。例如已知 1 至 40 与 43 的值，41、42 尚未决定，应用就只能执行到 40。直接先做 43，会让余额和编号依赖各副本自己的接收先后。新领导者为缺口执行恢复，继承其中已有的候选。对准备后确实可以自由选值的位置，可以选用空命令填补。等 41、42 都决定以后，状态机再依次执行到 43~\cite{paxos-made-simple}。《The Part-Time Parliament》也分别讨论了按位置形成命令序列、补齐空洞，以及用当前状态和最后已纳入位置代替无限增长的历史~\cite{part-time-parliament}。

\CElevenFigure{ch11-original17}{各位置分别决定可以暂留空洞，应用等待连续前缀，同一 S1 也可能位于不同位置}{fig:ch11-paxos-prefix}

假如 S1 在 41 已经选定并应用，客户端没有收到结果，又将 S1 送到新领导者，后者可以把同一命令放到 43。Paxos 对位置 41 和 43 分别取得唯一值，这条日志完全符合协议。状态机依然要读取 $M(S1)$，恢复原 P1/O1/D1 和原 7000。改用 Paxos 后，第~\ref{section:11-3}~节的身份绑定、原输入比较与原结果保存有同样的作用。

Raft 采用另一种接班安排：候选者要带着足够新的日志取得选票，合法新领导者因此拥有旧提交前缀。经典 Paxos 的新提议者可以通过准备阶段收集旧接受事实，补足自己原先缺少的内容。两种安排把状态转移工作放在不同位置，其安全性都要沿相应协议的继承规则证明。Raft 对当前任期项的多数提交规则，也由它自己的日志与选举约束推导。

\CElevenFigure{ch11-paxos-raft}{Raft 在候选日志资格中保住前缀，Paxos 通过准备答复恢复旧接受事实}{fig:ch11-paxos-raft}

Howard 与 Mortier 把一份明确的 Multi-Paxos 表述放进与 Raft 相近的日志模型，突出接班时状态转移放在哪里：Raft 在投票时筛选日志足够新的候选者，选举消息主要比较末项任期和位置；他们所比较的 Paxos 允许本地缺项者先争取领导权，再从准备答复中补齐旧接受值~\cite{paxos-vs-raft}。因此，本节 L2 本地没有 41，仍可从 B 的答复继承 S1；第~\ref{section:11-2}~节的 Raft 则把旧前缀的保留放进候选资格与投票规则。

Wang 等人进一步问，既然两类协议都保存旧决定，能否把一类优化按规则迁移？原 Raft 与 Paxos 存在后缀处理等细节差异，论文先构造变体 Raft*，再用细化映射把它的状态和步骤解释为 Paxos 的状态及步骤。自动迁移针对文中定义的 non-mutating 类优化：它们增加辅助状态或步骤，但不改写基础 Paxos 状态的含义。论文以 Mencius 的负载分配和 Paxos Quorum Lease 的读取优化为案例~\cite{paxos-raft-parallels}。迁移产物是协议规格，开发者随后实现并验证持久化和故障路径。把这项工作用于本书，先核对所选库究竟实现哪份规格，再把 W/M/$\Gamma$、快照和原购买结果接到它实际保证的提交边界。

\subsection*{主备复制怎样延续一份确认}

主备复制描述了由主副本受理、向备副本传播的组织方式。一次确认由谁参与、配置怎样更换、旧主怎样停止，要由具体协议规定。要判断本章 DB-A 失效时哪项购买会留下来，就要沿这些确认与交接规则追踪。

PacificA 提供一条具体路径。当前配置指定主副本和各个备副本，每项更新由主副本分配连续序号。各副本先按序加入准备列表，主副本在收到当前配置中全部副本的准备确认后，才把连续提交点推进，并向客户端确认。备副本可以通过主副本后来携带的提交点，获知哪些准备项已经提交~\cite{pacifica}。

设当前配置为主 A、备 B/C，提交前缀到 40。A 把 S1 放到 41，B/C 都准备到 41，A 推进提交点后确认客户端。B 的本地提交点暂时停在 40，它的准备列表却已经含有 41。PacificA 在当前配置内保持的关系是：任意副本 q 的提交前缀，包含于主 p 的提交前缀，主的提交前缀又包含于 q 的准备列表~\cite{pacifica}：
\[
  \mathrm{committed}_q\subseteq
  \mathrm{committed}_p\subseteq
  \mathrm{prepared}_q.
\]

这里的包含关系按连续前缀解释。它使接班者的准备列表足以覆盖旧主已经确认的提交，尽管接班者可能尚未获知旧提交点。新主开始服务以前，先在新配置中完成协调：让自己准备列表中的有关项完成提交，再使各副本的准备列表与这份提交前缀衔接。刚才的 B 可以由此保住 41，再应用 S1。额外的未提交候选按协调规则保留或截去~\cite{pacifica}。

\CElevenFigure{ch11-original18}{PacificA 接班从准备列表保住旧提交，配置管理器与租约另行约束主资格}{fig:ch11-pacifica}

配置管理器保存当前配置及版本，变更请求携带旧版本，匹配时才安装更高版本。A 想删掉失联的 B/C，B 又想提升自己，两份基于同一版本的竞争变更只会有一份先成功。与此同时，旧主本地可能还不知道版本已变。PacificA 的租约机制让旧主在本地资格到期时停止服务，备副本经过更长的宽限期才申请取代它。论文采用有界时钟漂移条件来协调这两个期限~\cite{pacifica}。这里的期限比较主副本资格，第八章的任务领取期限则控制一项恢复任务的执行权。

Vertical Paxos 从另一角度连接主备组织与共识。它把配置选择交给辅助配置主控，并区分恢复旧事实时使用的读仲裁集合和接受新值时使用的写仲裁集合。把读集合缩成任意一个旧接受者，相应写集合就要覆盖全部接受者，保证两者相交。若接班者从旧接受者中选择，它本身就能提供所需旧事实~\cite{vertical-paxos}。

这种安排允许一个数据组用 k+1 个接受者容忍 k 个数据节点失效，系统同时依赖高可靠配置主控。配置主控本身的复制与可用性成本要计入部署。Vertical Paxos II 在旧状态转移完成后才激活新配置，适合解释单一活动配置的主备系统。Vertical Paxos I 则让新配置先活动，保留对旧配置的访问，完成转移后再解除依赖~\cite{vertical-paxos}。两者分别规定状态转移与激活的先后，使接班者能够继承旧决定。

\CElevenFigure{ch11-config-transfer}{Vertical Paxos I 与 II 把配置激活放在不同位置，旧事实的转移各有完成条件}{fig:ch11-config-transfer}

回到购买，若上述复制单元只含余额，接班者就只能继承余额。若命令与快照同时包含 W/M/$\Gamma$，原购买及其可恢复解释才会一起延续。应用交给复制层的内容，决定复制安全最终保存了什么。数据库高可用则把这个边界放在数据库事务与其恢复日志中，第~\ref{section:11-1}~节的 A/R 继续通过共同提交建立关系。

换副本和删日志都要把业务内容、身份解释、原结果与覆盖边界共同延续。随着有效身份增多，M 中保存的结果也会增多。若 E1 已经关闭，怎样删除一部分结果，又让迟到的旧 S1 得到正确处理，同时让已受理的 C1/T1 完成恢复，便成为下一章的具体问题。
\FloatBarrier
